\documentclass[10pt,twocolumn]{article}

\usepackage{amsmath,amssymb}
\usepackage{graphicx}

\setlength{\textwidth}{7.05in}
\setlength{\textheight}{9.25in}
\setlength{\oddsidemargin}{-0.30in}
\setlength{\evensidemargin}{-0.30in}
\setlength{\topmargin}{-0.55in}
\setlength{\columnsep}{0.25in}
\setlength{\parskip}{2pt}
\setlength{\parindent}{10pt}
\setlength{\floatsep}{5pt}
\setlength{\textfloatsep}{7pt}
\setlength{\intextsep}{5pt}

\renewcommand{\abstractname}{Resumen}
\renewcommand{\figurename}{Figura}
\renewcommand{\tablename}{Tabla}
\renewcommand{\refname}{Referencias}

\title{\vspace{-0.2in}Residuo barionico e inversion de materia total:\\
un marco candidato conectado con materia oscura asimetrica}
\author{German Garcia, A.\\Investigador independiente}
\date{Agosto 2026}

\begin{document}
\maketitle

\begin{abstract}
Se presenta una version de dos columnas, preparada como manuscrito tipo arXiv, de un marco candidato que compara la lectura residual de la asimetria barionica con una inversion algebraica extendida a la materia total cosmologica. El objetivo no es demostrar el origen fisico de la asimetria materia-antimateria ni identificar la naturaleza de la materia oscura, sino ordenar una regularidad cuantitativa bajo terminos verificables. Con $\eta_B=6.120\times10^{-10}$, la inversion barionica cierra con el contador radiativo usado en cosmologia. Al declarar la materia total como residuo efectivo, el cierre deja de ser unitario y produce $6.39$, compatible con $\Omega_m/\Omega_b$ para parametros de Planck. La prueba por densidad reproduce el mismo patron: cierre barionico $1.001$ y cierre total $6.393$. El resultado se interpreta como una pista candidata conectable con modelos de materia oscura asimetrica, no como un mecanismo dinamico completo.
\end{abstract}

\noindent\textbf{Palabras clave:} asimetria barionica; materia oscura asimetrica; CMB; BBN; deuterio; longitud Compton.

\section{Planteamiento}

La abundancia barionica actual suele resumirse mediante
\begin{equation}
\eta_B=\frac{n_B}{n_\gamma}\simeq 6.1\times10^{-10},
\end{equation}
donde $n_B$ es la densidad numerica neta de bariones y $n_\gamma$ la densidad numerica de fotones del fondo cosmico. Este valor es consistente entre nucleosintesis primordial y CMB \cite{pdg_bbn_2022,planck2018params}. En lenguaje fisico, la materia barionica visible puede leerse como un residuo neto de una asimetria materia-antimateria. Las condiciones necesarias para generar una asimetria dinamica fueron formuladas por Sakharov \cite{sakharov1967}; el mecanismo microscopico preciso sigue abierto \cite{dine_kusenko_2003,garbrecht2020}.

Aqui se propone una lectura inversa. Primero se toma el residuo barionico $N_B$ y se reconstruye el sector cancelado que seria compatible con el contador de fotones. Despues se repite la misma operacion declarando la materia total cosmologica como residuo efectivo. Si la segunda inversion produce el cociente observado entre materia total y materia barionica, el resultado puede compararse con la motivacion de los modelos de materia oscura asimetrica \cite{petraki_volkas_2013,zurek2014,boucenna_morisi_2014}.

\noindent\textbf{Guia de lectura.} En este manuscrito ``residuo'' significa componente neto observado despues de una cancelacion algebraica. No se afirma que existan dominios macroscopicos actuales de antimateria, ni que el CMB quede explicado por este calculo. El termino ``sector cancelado'' nombra una reconstruccion numerica candidata, no una poblacion observada directamente.

\section{Escala triadica y control Compton}

El barrido de escala usa
\begin{equation}
L_n=\ell_P3^n,\qquad M_n=m_P3^{-n},
\end{equation}
de modo que el producto masa-longitud queda fijo:
\begin{equation}
M_nL_n=m_P\ell_P=\frac{\hbar}{c}.
\end{equation}
La posicion continua de una masa $m$ sobre esta escalera se calcula con
\begin{equation}
n_C(m)=\log_3\left(\frac{\bar{\lambda}_C}{\ell_P}\right)
=\log_3\left(\frac{m_P}{m}\right),
\qquad
\bar{\lambda}_C=\frac{\hbar}{mc}.
\end{equation}
Este barrido no es una dinamica de formacion de particulas. Es un localizador de escala: pregunta donde caen electron, proton, neutron y deuteron si se comparan sus longitudes Compton reducidas con $L_n$.

\begin{table}[h!]
\centering
\caption{Sectores Compton principales. La tabla muestra la longitud Compton reducida, el nivel continuo $n_C$, el entero cercano y el cociente $\bar{\lambda}_C/L_n$.}
\scriptsize
\begin{tabular}{lrrrr}
\hline
Objeto & $\bar{\lambda}_C$ (m) & $n_C$ & $n$ & ratio\\
\hline
electron & $3.862{\times}10^{-13}$ & 46.903 & 47 & 0.899\\
proton/$^1$H & $2.103{\times}10^{-16}$ & 40.062 & 40 & 1.070\\
neutron & $2.100{\times}10^{-16}$ & 40.061 & 40 & 1.069\\
deuteron/$^2$H & $1.052{\times}10^{-16}$ & 39.431 & 39 & 1.606\\
\hline
\end{tabular}
\end{table}

La lectura inmediata es simple: proton y neutron se agrupan cerca de $n\simeq40$, electron cerca de $n\simeq47$, y deuteron se desplaza por su mayor masa nuclear. Esta organizacion sirve como control de escala para el contraste posterior con hidrogeno y deuterio.

\section{Control hidrogeno-deuterio}

El cierre nuclear usado como contraste es
\begin{equation}
^1{\rm H}=p,\qquad ^2{\rm H}=D=(p+n)_{\rm ligado},
\end{equation}
\begin{equation}
p+n\rightarrow D+\gamma.
\end{equation}
Con masas CODATA/NIST \cite{nist_codata_2022},
\begin{align}
m_p+m_n &=3.347549426510\times10^{-27}{\rm\,kg},\\
m_D &=3.343583776800\times10^{-27}{\rm\,kg},
\end{align}
y $\Delta m=3.965649710000\times10^{-30}{\rm\,kg}$. Por tanto, el deuteron debe tratarse como estado ligado y no como suma libre de proton mas neutron.

Para deuterio primordial, Cooke et al. reportan
\begin{equation}
(D/H)_P=2.527\times10^{-5}\cite{cooke2018}.
\end{equation}
La regularidad candidata
\begin{equation}
\sqrt{\eta_B}=2.473863\times10^{-5}
\end{equation}
produce
\begin{equation}
\Xi_D=\frac{(D/H)_P}{\sqrt{\eta_B}}=1.02148.
\end{equation}
La diferencia relativa es $2.10\%$. Esta cercania no sustituye a BBN; se reporta como un contraste posterior que el marco debe explicar si se desarrolla como modelo dinamico.

\section{Inversion del residuo}

La inversion algebraica parte de
\begin{equation}
N_B=N_b^0-N_{\bar b}^0,\qquad
N_b^0=N_{\rm par}+N_B,\qquad
N_{\bar b}^0=N_{\rm par}.
\end{equation}
Para conectar el componente cancelado con el contador radiativo se introduce $k_\gamma$:
\begin{equation}
N_{\bar b}^0=\frac{N_\gamma}{k_\gamma}
=\frac{N_B}{k_\gamma\eta_B}.
\end{equation}
Con $N_\gamma=1.465\times10^{89}$ y $N_B=8.966\times10^{79}$ se obtiene la inversion barionica. Si en lugar de usar solo bariones se usa la materia total como residuo efectivo, aparece
\begin{equation}
\frac{\Omega_m}{\Omega_b}=6.388,
\end{equation}
con $\Omega_m=0.315$ y $\Omega_b=0.049309$ \cite{planck2018params}.

\begin{table}[h!]
\centering
\caption{Inversion por cantidad total. El cierre compara el sector reconstruido con el contador radiativo normalizado; el caso barionico queda cerca de uno y el caso de materia total recupera $\Omega_m/\Omega_b$.}
\scriptsize
\begin{tabular}{lrrr}
\hline
Residuo & inv. $k=1$ & inv. $k=2$ & cierre\\
\hline
bariones & $1.47{\times}10^{89}$ & $7.33{\times}10^{88}$ & 1.000\\
materia total & $9.36{\times}10^{89}$ & $4.68{\times}10^{89}$ & 6.388\\
\hline
\end{tabular}
\end{table}

\section{Prueba por densidad}

La misma prueba puede hacerse sin escoger un volumen cosmologico. Con $H_0=67.4$ km s$^{-1}$ Mpc$^{-1}$,
\begin{equation}
\rho_{\rm crit}=\frac{3H_0^2}{8\pi G}=8.533\times10^{-27}{\rm\,kg\,m^{-3}}.
\end{equation}
Por tanto
\begin{align}
\rho_b&=4.207\times10^{-28}{\rm\,kg\,m^{-3}},\\
\rho_m&=2.688\times10^{-27}{\rm\,kg\,m^{-3}}.
\end{align}
Con $T_{\rm CMB}=2.72548$ K \cite{fixsen2009},
\begin{align}
n_\gamma&=4.107\times10^8{\rm\,m^{-3}},\\
n_b&=\rho_b/m_p=2.516\times10^{-1}{\rm\,m^{-3}},
\end{align}
de donde
\begin{equation}
\eta_{\rm densidad}=6.125\times10^{-10}.
\end{equation}

\begin{table}[h!]
\centering
\caption{Inversion por densidad. La prueba evita fijar el radio del universo observable y conserva el mismo patron: cierre barionico unitario y cierre total cercano a $6.39$.}
\scriptsize
\begin{tabular}{lrrr}
\hline
Residuo & $\rho_{\rm canc}$ $k=1$ & $\rho_{\rm canc}$ $k=2$ & cierre\\
\hline
$\rho_b$ & $6.875{\times}10^{-19}$ & $3.437{\times}10^{-19}$ & 1.001\\
$\rho_m$ & $4.392{\times}10^{-18}$ & $2.196{\times}10^{-18}$ & 6.393\\
\hline
\end{tabular}
\end{table}

\begin{figure*}[t]
\centering
\includegraphics[width=0.88\textwidth]{arxiv_fig_residuo_cierres_bw_v3.png}
\caption{Comparacion visual de la inversion residual. El panel A muestra, en escala logaritmica, bariones netos, materia total equivalente, fotones CMB usados como contador y sectores cancelados reconstruidos para $k_\gamma=1$. El panel B resume los cierres normalizados: los casos barionicos quedan cerca de uno, mientras que los casos de materia total reproducen el factor cosmologico cercano a $6.39$. La figura usa escala de grises y patrones para no depender del color.}
\end{figure*}

\section{Materia oscura asimetrica}

La literatura de materia oscura asimetrica parte de una motivacion cercana: la densidad de materia oscura es mayor que la densidad barionica, pero no por decenas de ordenes de magnitud \cite{petraki_volkas_2013,zurek2014,boucenna_morisi_2014}. Si la asimetria numerica oscura fuera comparable a la barionica,
\begin{equation}
m_{\rm DM}\sim \frac{\Omega_{\rm DM}}{\Omega_b}m_p.
\end{equation}
Usando los valores anteriores,
\begin{equation}
\frac{\Omega_{\rm DM}}{\Omega_b}
=\frac{\Omega_m-\Omega_b}{\Omega_b}=5.388,
\end{equation}
lo que da una escala candidata
\begin{equation}
m_{\rm DM}\sim5.05{\rm\,GeV}.
\end{equation}
Este valor no es una prediccion unica del presente marco. Es la escala familiar que aparece cuando se intenta relacionar una asimetria visible con una asimetria oscura.

\begin{figure*}[t]
\centering
\includegraphics[width=0.86\textwidth]{arxiv_fig_compton_deuterio_bw_v3.png}
\caption{Barrido triadico Planck-Compton y contraste posterior con deuterio. El panel A ubica electron, proton, neutron y deuteron mediante $n_C=\log_3(\bar{\lambda}_C/\ell_P)$ sobre la escalera $L_n=\ell_P3^n$. El panel B compara la abundancia primordial observada $(D/H)_P$ con $\sqrt{\eta_B}$ en unidades de $10^{-5}$. El resultado es una cercania numerica candidata, no una sustitucion de los calculos de nucleosintesis primordial.}
\end{figure*}

\section{Discusion y limites}

El resultado conservador es el siguiente. Primero, la lectura barionica como residuo esta alineada con la asimetria barionica estandar. Segundo, la inversion barionica cierra con el contador de fotones por la definicion de $\eta_B$. Tercero, al declarar materia total como residuo efectivo aparece $\Omega_m/\Omega_b\simeq6.39$. Cuarto, la prueba por densidad reproduce el mismo factor sin depender del radio del universo observable. Quinto, el puente con materia oscura asimetrica es razonable como programa teorico, pero no esta demostrado por este calculo.

Las condiciones pendientes son sustanciales: falta una lagrangiana o red cinetica, tasas de aniquilacion, transferencia de asimetria entre sectores, compatibilidad completa con BBN/CMB y limites de deteccion directa. Ademas, dominios macroscopicos de materia-antimateria estan fuertemente restringidos por rayos gamma y por el CMB \cite{cohen_derujula_glashow_1998}.

\section{Conclusion}

La inversion del residuo barionico reproduce el cierre radiativo esperado por $\eta_B$. Al aplicar la misma inversion a la materia total, el cierre obtenido es $6.39$, coincidente con $\Omega_m/\Omega_b$. Esta coincidencia convierte el calculo en una regularidad candidata conectable con materia oscura asimetrica. El siguiente paso es dinamico: reproducir $\eta_B$, $\Omega_{\rm DM}/\Omega_b$, $(D/H)_P$ y las restricciones CMB/BBN.

\section*{Declaracion de alcance}

Este manuscrito no afirma que el origen de la asimetria materia-antimateria ni la naturaleza de la materia oscura queden demostrados. Presenta un marco candidato de consistencia cuantitativa, formulado para que sus entradas, salidas y limites puedan revisarse de manera independiente.

\begin{thebibliography}{99}
\bibitem{nist_codata_2022} P. J. Mohr, E. Tiesinga, D. B. Newell, and B. N. Taylor. CODATA 2022 Values of the Fundamental Physical Constants. NIST, 2024.
\bibitem{planck2018params} Planck Collaboration. Planck 2018 results. VI. Cosmological parameters. Astron. Astrophys. 641:A6, 2020. doi:10.1051/0004-6361/201833910.
\bibitem{pdg_bbn_2022} Particle Data Group. Big-Bang Nucleosynthesis review. Review of Particle Physics, 2022.
\bibitem{cooke2018} R. J. Cooke, M. Pettini, and C. C. Steidel. One Percent Determination of the Primordial Deuterium Abundance. Astrophys. J. 855:102, 2018. doi:10.3847/1538-4357/aaab53.
\bibitem{fixsen2009} D. J. Fixsen. The Temperature of the Cosmic Microwave Background. Astrophys. J. 707:916--920, 2009. doi:10.1088/0004-637X/707/2/916.
\bibitem{sakharov1967} A. D. Sakharov. Violation of CP Invariance, C Asymmetry, and Baryon Asymmetry of the Universe. JETP Lett. 5:24--27, 1967.
\bibitem{dine_kusenko_2003} M. Dine and A. Kusenko. Origin of the matter-antimatter asymmetry. Rev. Mod. Phys. 76:1--30, 2003. doi:10.1103/RevModPhys.76.1.
\bibitem{garbrecht2020} B. Garbrecht. Why is there more matter than antimatter? Prog. Part. Nucl. Phys. 110:103727, 2020. doi:10.1016/j.ppnp.2019.103727.
\bibitem{cohen_derujula_glashow_1998} A. G. Cohen, A. De Rujula, and S. L. Glashow. A Matter-Antimatter Universe? Astrophys. J. 495:539--549, 1998. doi:10.1086/305328.
\bibitem{petraki_volkas_2013} K. Petraki and R. R. Volkas. Review of asymmetric dark matter. Int. J. Mod. Phys. A 28:1330028, 2013. doi:10.1142/S0217751X13300287.
\bibitem{zurek2014} K. M. Zurek. Asymmetric Dark Matter: Theories, Signatures, and Constraints. Phys. Rep. 537:91--121, 2014. doi:10.1016/j.physrep.2013.12.001.
\bibitem{boucenna_morisi_2014} S. M. Boucenna and S. Morisi. Theories relating baryon asymmetry and dark matter. Front. Phys. 1:33, 2014. doi:10.3389/fphy.2013.00033.
\end{thebibliography}

\end{document}
